\summarychapter{13}{The Sage of the Ages}{Augustine}

\begin{inbrief}
When Rome fell to the Visigoths in 410, Augustine, bishop of Hippo in North Africa, set out to explain
it. The chapter tells of his long road to conversion, his answers to the Donatists on the church and
to Pelagius on sin and grace, and his vision of two cities, which gave Christendom its philosophy.
Roman Catholics still draw on his doctrine of the church, Protestants on his teaching on sin and grace.
\question{Why had Rome fallen? Would the ruin of the Eternal City mean the collapse of Christianity, or
the end of the world?}
\end{inbrief}

For 620 years, since Hannibal, no foreign enemy had stood before Rome. In 410 Alaric and his Arian
Visigoths besieged and plundered it, sparing only the churches, to which Alaric had the sacred vessels
carried. Rome's glory was gone. ``The city that took the whole world captive is itself taken,'' said
Jerome. Pagans asked whether the old gods were angry at being abandoned for the Christian God.
Watching the refugees arrive at Hippo, Augustine decided to answer.

\ssection{Preparation for Greatness}
Augustine was born on November 13, 354, at Tagaste in Numidia (today's Algeria), son of an easygoing
pagan, Patricius, and an ardent Christian, Monica. Schooled at home and at Carthage, he gave way to
sexual temptation and lived thirteen years with a woman who bore him a son, Adeodatus; he always felt
sex to be his defiling passion, and it colored his idea of sin. At nineteen a treatise of Cicero
persuaded him to seek truth, but the Bible seemed crude to him. For nine years (374--383) he was a
Manichean. Mani, crucified in Persia in 276 or 277, had taught an eternal conflict of good and evil in
which man's spirit is good and his body evil; the good is freed by prayer and abstinence. Like the
Gnostics, Manicheans denied that Jesus had a real body or really died and rejected the Old Testament.
Augustine taught grammar at Tagaste and rhetoric at Carthage, moved to Rome in 383 and in 384 became
professor at Milan. At thirty, at the height of success, he was deeply unhappy: he sent away his
mistress to become engaged to an heiress and still could not master his passions.

In Milan the preaching of Bishop Ambrose showed him that Christianity could be eloquent and
intelligent and that hard Old Testament stories could be read as allegories. The story of Anthony and
the Egyptian hermits shamed him: unlearned men conquered where he, with all his education, failed. In
agony in his garden he heard a child's voice chant, ``Take it and read it,'' opened the New Testament
at Romans 13:13--14, and felt ``the light of peace'' pour into his heart. Ambrose baptized him, with
Adeodatus and his friend Alypius, at Easter 387. On the way home his mother died near Rome, and in 388
his son died at Tagaste. He wished to withdraw from the world, but at Hippo he was ordained priest by
popular demand, then made assistant to Bishop Valerius and soon his successor. From the age of
forty-three until his death in 430 he stood at the center of the storms of his time.

\ssection{The Donatist Controversy}
North African Christianity had been split for almost a century. The Donatists, named after Donatus,
bishop of Carthage (313--355), stood for a holy church, strict discipline and resistance to unworthy
bishops. They charged that Catholic bishops had handed over the Scriptures to be burned under
Diocletian, an act of apostasy, so that clergy ordained by them were no true clergy and the Donatists
alone were the true church. In some districts they were the majority. Augustine rejected the idea of
a pure church: until the judgment the church is a mixed multitude of good and bad, as in the parable
of the wheat and the tares (though that parable, Shelley notes, speaks of the world, not the church).
He also denied that a sacrament's validity depends on the minister's morals. The sacrament belongs to
Christ, and the acts of a properly ordained priest are God's acts. This made the priest the channel of
grace and lent Augustine's weight to the priestly (sacerdotal) view of the church that went to
unfortunate extremes in the Middle Ages. At first opposed to coercion, he came step by step to accept
the use of force against the Donatists, citing ``Compel people to come in''; his prestige later served
those who justified the Inquisition.

\ssection{Of Sin and Grace}
The Pelagian debate brought out his teaching on grace. Pelagius, a British monk who came from Rome,
and his disciple Coelestius found no welcome in Augustine's Africa and went east, where churches were
more receptive. Pelagius denied that sin is inherited from Adam. Man is free to do right or wrong;
Adam harmed his descendants only by bad example; death is not the penalty of his sin; it is possible
not to sin, and some have lived without it. God predestines only in the sense of foreseeing who will
believe. Forgiveness comes by ``faith alone,'' but the forgiven have power in themselves to please
God, so the special help of the Spirit is hardly needed. It was practically the Stoic ideal of
self-control.

Augustine's experience said the opposite: only irresistible divine power had rescued him, and only
constantly renewed grace kept him. In Adam the whole race fell and became a ``mass of corruption''
incapable of any saving good and deserving damnation. All power to do good is God's free gift, grace.
From the fallen mass God chooses some to receive it, through Christ's work and ordinarily through the
church and its sacraments; baptism gives regenerating grace. The rest are lost, and no one can be sure
of salvation, since only those given the further grace of perseverance are finally redeemed. Salvation
is wholly of God. Augustine launched a literary campaign; the emperor Honorius banished the Pelagians
by 419, and the Council of Ephesus condemned them in 431. The Reformers later built on Augustine's
view of human helplessness, grace and predestination, while setting aside his tie between grace and
the sacraments. In the following century, however, the church kept its distance from irresistible
grace and predestination; critics called them an innovation against the long-held belief in human
freedom, and Vincent of Lerins insisted that Christians should believe what has been believed
``always, everywhere, and by all.''

\ssection{The City of God on the Ruins of Rome}
Augustine was fifty-six when Rome was sacked. He housed the refugees and preached that cities are men,
not walls: Rome, unlike Sodom, was chastised but not destroyed. Then for sixteen years he worked on
the deeper question, the relation between earthly cities that rise and fall and the everlasting City
of God. The result, \emph{The City of God}, shaped Christian thinking about God and Caesar for fifteen
centuries. From Adam to the end, humanity forms two cities: the godless, united by love of temporal
things, and the spiritual, born of grace and bound together by the love of God. The Romans' great
deeds were driven by love of glory, a kind of life after death on the lips of men; the heavenly city
outshines Rome, offering truth instead of victory, holiness instead of rank, eternity instead of life.
The church is the only human community that builds the City of God. The state has its place in
restraining crime and keeping peace, but since it rests on the power of sin it must submit to the laws
of the church. The book gave a spiritual meaning to the miseries of the age: the golden age lies ahead
in God's kingdom, not behind in a crumbling empire.

Augustine's last years were hard. In his seventy-sixth year the Vandals crossed from Spain and swept
toward Hippo. He had the penitential psalms fixed to the wall beside his bed and believed the end of
the world near. He died on August 28, 430. The siege lasted fourteen months; when the Vandals broke in
in August 431 they found most of the people dead or dying of hunger.

\begin{keyfacts}
\keypeople{Augustine; Monica; Patricius; Adeodatus; Alypius; Ambrose; Mani; Donatus; Pelagius;
Coelestius; Honorius; Vincent of Lerins; Alaric the Visigoth; Jerome.}
\keydates{%
  \kd{354}{Birth of Augustine at Tagaste (November 13)}
  \kd{374--383}{Augustine a Manichean}
  \kd{387}{Baptized by Ambrose at Milan}
  \kd{410}{Alaric and the Visigoths sack Rome}
  \kd{419}{The emperor Honorius banishes the Pelagians}
  \kd{430}{Death of Augustine (August 28) during the Vandal siege of Hippo}
  \kd{431}{Council of Ephesus condemns Pelagianism; Hippo falls}}
\keyterms{Manicheanism; Donatism; mixed church; validity of the sacraments; sacerdotal view; coercion;
Pelagianism; original sin; grace; predestination; perseverance; \emph{The City of God}.}
\end{keyfacts}
